% Nouveau EGA V (cote 25): archive cote 25, site page 147 (PDF page 148).
% typescript. Draft compared: none.
% AI-transcribed by Claude Opus 5.5 (claude-opus-5-5, adaptive thinking, effort high), Message Batches API, batch ega25b (msgbatch_01A7n5YuXubUPNK1zsepi8RG), 2026-11-13; not yet checked by eye.
\scanpage{25/147}
\mspage{8}

\uline{de} $\overset{\hand{$\vee$}}{\underline{P}}_{X/Y}\otimes\underline{O}_{V^1}$ , \uline{et en tant qu'homomorphisme dans c ette composante,}
\uline{il est déterminé précisément par lasedtion} $D = \text{dét}\,M(\phi)$ \uline{déjà envisagée}
\uline{par ailleurs} (tjrs modulo un signe à déterminer). Bien entendu,
l'homomorphisme b) s'en déduit, compte tenu de la commutativité énoncée
dans le lemme. \deleted{Aux}\inserted{Les} points non supersinguliers$_{\text{x}}$ sont donc exacte\unsure{m}ent
ceux en lesquels $\underline{a}$: $\underline{\omega}^{-2}_{X/Y}\otimes\underline{O}_{V^1}\longrightarrow\underline{O}_{\unsure{V}^1}$ est un isomorphisme, et en
ces points, \xout{\unsure{xxxxxxdéfinitxxxxisomorphismexxdu}} le facteur $\underline{O}_{\unsure{V}^1}$ s'identif
fie donc exactement à l'intersection de $\underline{P}^1_{X/Y}\otimes\underline{O}_{V^1}$ et $\underline{\Omega}^1_{Y/S}\otimes\underline{O}_{V^1}$
dans $\underline{\Omega}^1_{X/S}\otimes\underline{O}_{V^1}$ . Donc sur l'ouvert de\xout{x} $V^1$ formé des points non
supersinguliers, on a une suite exacte
\[ 0\longrightarrow\underline{O}_{V^1}\longrightarrow\underline{\Omega}^1_{Y/S}\otimes\underline{O}_{V^1}\longrightarrow\underline{\Omega}^1_{V^1/S}\longrightarrow 0 \qquad \xout{x}, \]
ainsi que la suite exacte
\[ 0\longrightarrow\underline{O}_{V^1}\longrightarrow\underline{\Omega}^1_{X/S}\otimes\underline{O}_{V^1}\longrightarrow\xout{x}(\underline{\Omega}^1_{X/Y}+\underline{\Omega}^1_{Y/S})\otimes\underline{O}_{V^1}\longrightarrow 0 \]

\margin{|} Il faudrait résumer  toute cette situation dans une proposition
récapitulative, et sous une forme légèrement plus générale, en partant
d'une section $\phi$ d'un \uline{Module inversible} \xout{xxx} K \xout{sur} X - ce qui oblige
à tordre par K un certain\xout{s} nombre\xout{s} des Modules introduits dans les
considérations précéde\unsure{nt|mm}es. C'est en effet  cette situation qu'on
rencontre  dans le par. des sections hyperplanes. Signaler aussi que
\margin{|} $V(\phi)^{\text{sing}}$ \xout{xxxxx} et $V(\phi)^{\text{sup sing}}$ ne changent pas en r emplaçant $\phi$ par u $\phi$,
où u est une unité, \hand{en fait $d\phi|V(\phi)$ et $M(\phi)$ sont multipliés par \unsure{$u$},}
\inserted{et $D(\phi)$ est multiplié par $u^n$ ($n$ = dim. relative de $X$ sur $Y$),}
\uline{Remarque} Sous les conditions envisagées sur X,Y,S,$\phi$ , comme \xout{\unsure{x(x}})
$U = V(\phi)^{\text{sing}}- V(\phi)^{\text{sup sing}}$ est non ramffié sur Y, ses fibres sur Y sont
discrètes, ce qui implique que pour tout $y\in Y$, tout zéro singulier ordi-%
naire de $\phi_y$ dans $X_y$ est isolé dans l'ensemble des zéros singuliers, \deleted{et}
\deleted{à fortiori est}\inserted{i.e. est} isolé dans l'ensemble des points non lisses de $V(\phi_y)/k(y)$.
\xout{Parxsuite} On en conclut \inserted{p.ex.} aussitôt que si $d\geqslant 2$, c'est un point géométri-%

\margin{147}
